\exercisesfor{3}

字母 X 表示一个集合。

\begin{enumerate}
\def\labelenumi{\arabic{enumi}.}
\tightlist
\item
  \begin{enumerate}
  \def\labelenumii{(\alph{enumii})}
  \tightlist
  \item
    证明每个元素 \(u \in \mathrm{A}^{+}(\mathrm{X})\) 都能唯一地写成 \(u = \sum_{x \in \mathrm{X}} u_{x} x\) 的形式，其中 \(u_{x} \in \mathrm{A}(\mathrm{X})\)。
  \end{enumerate}
\end{enumerate}

\begin{enumerate}
\def\labelenumi{(\alph{enumi})}
\setcounter{enumi}{1}
\tightlist
\item
  证明 \(\{0\}\) 是 \(\mathrm{A}^{+}(\mathrm{X})\) 中唯一在所有映射 \(u \mapsto u_{x} (x \in \mathrm{X})\) 下稳定的子模。（若 a 是这样的子模且 \(a \neq \{0\}\)，考虑 a 中次数最小的非零元素。）
\end{enumerate}

\begin{enumerate}
\def\labelenumi{\arabic{enumi}.}
\setcounter{enumi}{1}
\tightlist
\item
  设 g 是一个自由模李代数；通过 \(\sigma: g \to Ug\)，将其与它在包络代数 \(Ug\) 中的像相识别。令 \(U^{+}g\) 为典范同态 \(Ug \to K\) 的核。令 i 是从 X 到 g 的映射，使得 \(i(X)\) 作为李代数生成 g。
\end{enumerate}

\begin{enumerate}
\def\labelenumi{(\alph{enumi})}
\item
  证明 \(\mathbf{U}^{+}\mathfrak{g}\) 作为左 Ug-模由 \(i(\mathbf{X})\) 生成。
\item
  证明以下性质等价：
\item
  g 是以 i: X → g 为基本族的自由李代数。
\end{enumerate}

\begin{enumerate}
\def\labelenumi{(\roman{enumi})}
\setcounter{enumi}{1}
\tightlist
\item
  族 \(i\) 是左 Ug-模 \(\mathbf{U}^{+}\mathfrak{g}\) 的一组基。
\end{enumerate}

（蕴含关系 (i) \(\Rightarrow\) (ii) 利用同构 \(\mathrm{Ug} \to \mathrm{A}(\mathrm{X})\)，由习题 1 (a) 得出。为证明 (ii) \(\Rightarrow\) (i)，将习题 1 (b) 应用于同态 \(\mathrm{A}(\mathrm{X}) \to \mathrm{Ug}\) 的核，该同态由 \(i\) 定义。）

\begin{enumerate}
\def\labelenumi{\arabic{enumi}.}
\setcounter{enumi}{2}
\tightlist
\item
  设 \(x \in \mathbf{X}\)。证明 \(x\) 在 \(\mathrm{A}(\mathbf{X})\) 中的中心化子是由 \(x\) 生成的子代数。由此推出，\(\mathrm{L}(\mathbf{X})\) 中与 \(x\) 对易的唯一元素是 \(x\) 的倍数。特别地，\(\mathrm{L}(\mathbf{X})\) 的中心在 \(\operatorname{Card}(\mathbf{X}) \geqslant 2\) 时为 (0)，而 \(\mathrm{L}(\mathbf{X}) / \left( \sum_{n > p} \mathrm{L}^n(\mathbf{X}) \right)\) 的中心化为 \(\mathrm{L}^p(\mathbf{X})\) 的典范像。
\end{enumerate}


¶ 4. 设 K 是特征为 p \textgreater{} 0 的域。

\begin{enumerate}
\def\labelenumi{(\alph{enumi})}
\item
  设 \(\mathbf{H}\) 是 \(\mathbf{L}(\mathbf{X})\)（它是 \(\mathbf{A}(\mathbf{X})\) 的李子代数）的一组基。证明（利用同构 \(\mathbf{U}(\mathbf{L}(\mathbf{X})) \to \mathbf{A}(\mathbf{X})\) 和庞加莱–伯克霍夫–维特定理），元素 \(h^{p^n}\)（\(h \in \mathbf{H}\)、\(n \in \mathbf{N}\)）在 \(\mathbf{K}\) 上线性无关。
\item
  若 \(m\) 是整数 \(\geqslant 0\)，令 \(L_m\) 表示 \(A(X)\) 中以 \(h^{p^n}\)（其中 \(h \in H, n \leqslant m\)）为基的子模。证明 \(L_m\) 是 \(A(X)\) 的李子代数，并且若 \(a \in L_m - 1\)，则 \(a^p \in L_m\)（对 \(m\) 归纳，并使用 Jacobson 公式，参见~第一章，§ 1，习题 19）。推出 \(L_m\) 与 \(H\) 的选择无关，并且它是 \(A(X)\) 中由 \(x^{p^n}\)（其中 \(x \in X, n \leqslant m\)）生成的李子代数。
\item
  令 \(\mathbf{L}(\mathbf{X}, p)\) 为各 \(\mathbf{L}_m\)（\(m \geqslant 0\)）的并。证明 \(\mathbf{L}(\mathbf{X}, p)\) 是最小李 \(p\)-子代数，包含于 \(\mathbf{A}(\mathbf{X})\) 且包含 \(\mathbf{X}\)（参见~第一章 §1 的习题 20）；它以 \(h^{p^n}\)（其中 \(h \in \mathbf{H}\)、\(n \in \mathbf{N}\)）为基。
\item
  令 \(\tilde{\mathbf{U}}(\mathbf{L}(\mathbf{X},p))\) 为 \(\mathbf{L}(\mathbf{X},p)\) 的受限包络代数，参见~第一章 §2 的习题 6。证明 \(\mathbf{L}(\mathbf{X},p)\) 到 \(\mathbf{A}(\mathbf{X})\) 的嵌入可以延拓为 \(\tilde{\mathbf{U}}(\mathbf{L}(\mathbf{X},p))\) 到 \(\mathbf{A}(\mathbf{X})\) 的同构。（使用庞加莱–伯克霍夫–维特定理和上面引用的习题。）推出 \(\mathbf{L}(\mathbf{X},p)\) 是 \(\mathbf{A}(\mathbf{X})\) 的本原元素集合，参见~§1 的习题 12。
\item
  设 \(f\) 是从 \(\mathbf{X}\) 到李 \(p\)-代数 \(\mathfrak{g}\) 的映射。证明 \(f\) 可以唯一延拓为 \(p\)-同态 \(\mathbf{F} \colon \mathbf{L}(\mathbf{X}, p) \to \mathfrak{g}\)。（先将 \(f\) 延拓为从 \(\mathrm{A}(\mathbf{X})\) 到 \(\tilde{\mathrm{U}}\mathfrak{g}\) 的代数同态。）
\end{enumerate}

（李 \(p\)-代数 \(\mathbf{L}(\mathbf{X}, p)\) 称为自由李 \(p\)-代数（集合 \(\mathbf{X}\) 上）。）

